\documentclass[10pt]{article}
\usepackage[utf8]{inputenc}
\usepackage[T1]{fontenc}
\usepackage{mathptmx}
\usepackage[margin=1in]{geometry}
\usepackage[numbers,sort&compress]{natbib}
\usepackage{graphicx}
\usepackage{amsmath,amssymb}
\usepackage{booktabs}
\usepackage{xcolor}
\usepackage[font=small,labelfont=bf]{caption}
\usepackage[hidelinks]{hyperref}

\newcommand{\figifexists}[2][]{\IfFileExists{#2}{\includegraphics[#1]{#2}}{\fbox{\parbox{2in}{\centering\small Missing:\\\texttt{\detokenize{#2}}}}}}
% AUTO-GENERATED by src/paper_numbers.py -- do not edit by hand
\newcommand{\NOSS}{105}
\newcommand{\AccDirectOSS}{45.7}
\newcommand{\TokDirectOSS}{996}
\newcommand{\AccVoteOSS}{47.6}
\newcommand{\TokVoteOSS}{2,909}
\newcommand{\AccCritiqueOSS}{44.8}
\newcommand{\TokCritiqueOSS}{3,227}
\newcommand{\AccDebateOSS}{48.6}
\newcommand{\TokDebateOSS}{3,211}
\newcommand{\DDebateMinusVoteOSS}{+1.0}
\newcommand{\DDebateMinusVoteOSSLo}{-7.6}
\newcommand{\DDebateMinusVoteOSSHi}{+9.5}
\newcommand{\PDebateMinusVoteOSS}{1.00}
\newcommand{\DCritiqueMinusVoteOSS}{-2.9}
\newcommand{\DCritiqueMinusVoteOSSLo}{-14.3}
\newcommand{\DCritiqueMinusVoteOSSHi}{+8.6}
\newcommand{\PCritiqueMinusVoteOSS}{0.74}
\newcommand{\DVoteMinusDirectOSS}{+1.9}
\newcommand{\DVoteMinusDirectOSSLo}{-4.8}
\newcommand{\DVoteMinusDirectOSSHi}{+8.6}
\newcommand{\PVoteMinusDirectOSS}{0.77}
\newcommand{\DisFracOSS}{37}
\newcommand{\DisAccDirectOSS}{31}
\newcommand{\DisAccVoteOSS}{31}
\newcommand{\DisAccCritiqueOSS}{33}
\newcommand{\DisAccDebateOSS}{33}
\newcommand{\RescueCritiqueOSS}{31}
\newcommand{\CorruptCritiqueOSS}{40}
\newcommand{\RescueDebateOSS}{20}
\newcommand{\CorruptDebateOSS}{20}
\newcommand{\CritMinusVoteSeqQAOSS}{-24}
\newcommand{\CritMinusVoteProtocolOSS}{-5}
\newcommand{\CritMinusVoteSeqQAIIOSS}{-10}
\newcommand{\NGemma}{127}
\newcommand{\AccDirectGemma}{48.8}
\newcommand{\TokDirectGemma}{1,091}
\newcommand{\AccVoteGemma}{49.6}
\newcommand{\TokVoteGemma}{3,233}
\newcommand{\AccCritiqueGemma}{45.7}
\newcommand{\TokCritiqueGemma}{3,766}
\newcommand{\AccDebateGemma}{56.7}
\newcommand{\TokDebateGemma}{4,172}
\newcommand{\DDebateMinusVoteGemma}{+7.1}
\newcommand{\DDebateMinusVoteGemmaLo}{+0.0}
\newcommand{\DDebateMinusVoteGemmaHi}{+14.2}
\newcommand{\PDebateMinusVoteGemma}{0.08}
\newcommand{\DCritiqueMinusVoteGemma}{-3.9}
\newcommand{\DCritiqueMinusVoteGemmaLo}{-10.3}
\newcommand{\DCritiqueMinusVoteGemmaHi}{+2.4}
\newcommand{\PCritiqueMinusVoteGemma}{0.36}
\newcommand{\DVoteMinusDirectGemma}{+0.8}
\newcommand{\DVoteMinusDirectGemmaLo}{-5.5}
\newcommand{\DVoteMinusDirectGemmaHi}{+6.3}
\newcommand{\PVoteMinusDirectGemma}{1.00}
\newcommand{\DisFracGemma}{43}
\newcommand{\DisAccDirectGemma}{31}
\newcommand{\DisAccVoteGemma}{31}
\newcommand{\DisAccCritiqueGemma}{24}
\newcommand{\DisAccDebateGemma}{46}
\newcommand{\RescueCritiqueGemma}{11}
\newcommand{\CorruptCritiqueGemma}{19}
\newcommand{\RescueDebateGemma}{23}
\newcommand{\CorruptDebateGemma}{10}
\newcommand{\CritMinusVoteSeqQAGemma}{+0}
\newcommand{\CritMinusVoteProtocolGemma}{+4}
\newcommand{\CritMinusVoteSeqQAIIGemma}{-12}
\newcommand{\NLite}{75}
\newcommand{\AccDirectLite}{61.3}
\newcommand{\TokDirectLite}{1,134}
\newcommand{\AccVoteLite}{62.7}
\newcommand{\TokVoteLite}{3,143}
\newcommand{\AccCritiqueLite}{56.0}
\newcommand{\TokCritiqueLite}{3,179}
\newcommand{\AccDebateLite}{57.3}
\newcommand{\TokDebateLite}{3,609}
\newcommand{\DDebateMinusVoteLite}{-5.3}
\newcommand{\DDebateMinusVoteLiteLo}{-13.3}
\newcommand{\DDebateMinusVoteLiteHi}{+1.3}
\newcommand{\PDebateMinusVoteLite}{0.29}
\newcommand{\DCritiqueMinusVoteLite}{-6.7}
\newcommand{\DCritiqueMinusVoteLiteLo}{-17.3}
\newcommand{\DCritiqueMinusVoteLiteHi}{+4.0}
\newcommand{\PCritiqueMinusVoteLite}{0.36}
\newcommand{\DVoteMinusDirectLite}{+1.3}
\newcommand{\DVoteMinusDirectLiteLo}{-4.0}
\newcommand{\DVoteMinusDirectLiteHi}{+6.7}
\newcommand{\PVoteMinusDirectLite}{1.00}
\newcommand{\DisFracLite}{27}
\newcommand{\DisAccDirectLite}{30}
\newcommand{\DisAccVoteLite}{35}
\newcommand{\DisAccCritiqueLite}{25}
\newcommand{\DisAccDebateLite}{35}
\newcommand{\RescueCritiqueLite}{25}
\newcommand{\CorruptCritiqueLite}{26}
\newcommand{\RescueDebateLite}{7}
\newcommand{\CorruptDebateLite}{13}
\newcommand{\CritMinusVoteSeqQALite}{-20}
\newcommand{\CritMinusVoteProtocolLite}{+7}
\newcommand{\CritMinusVoteSeqQAIILite}{-20}
\newcommand{\TDirectLowLite}{70}
\newcommand{\TTokDirectLowLite}{478}
\newcommand{\TDirectHighLite}{67}
\newcommand{\TTokDirectHighLite}{2,539}
\newcommand{\TVoteLowLite}{70}
\newcommand{\TTokVoteLowLite}{1,141}
\newcommand{\TCritiqueLowLite}{63}
\newcommand{\TTokCritiqueLowLite}{1,004}
\newcommand{\TDebateLowLite}{67}
\newcommand{\TTokDebateLowLite}{1,440}
\newcommand{\TDirectLowLuna}{47}
\newcommand{\TTokDirectLowLuna}{115}
\newcommand{\TDirectHighLuna}{60}
\newcommand{\TTokDirectHighLuna}{1,839}
\newcommand{\TVoteLowLuna}{53}
\newcommand{\TTokVoteLowLuna}{341}
\newcommand{\TVoteHighLuna}{63}
\newcommand{\TTokVoteHighLuna}{6,016}
\newcommand{\TCritiqueLowLuna}{53}
\newcommand{\TTokCritiqueLowLuna}{217}
\newcommand{\TCritiqueHighLuna}{63}
\newcommand{\TTokCritiqueHighLuna}{5,267}
\newcommand{\TDebateLowLuna}{60}
\newcommand{\TTokDebateLowLuna}{280}
\newcommand{\TDebateHighLuna}{63}
\newcommand{\TTokDebateHighLuna}{5,809}
\newcommand{\TDirectHighSol}{70}
\newcommand{\TTokDirectHighSol}{3,228}
\newcommand{\TVoteHighSol}{70}
\newcommand{\TTokVoteHighSol}{8,753}
\newcommand{\TDirectLowOSS}{50}
\newcommand{\TTokDirectLowOSS}{284}
\newcommand{\TVoteLowOSS}{50}
\newcommand{\TTokVoteLowOSS}{836}
\newcommand{\TCritiqueLowOSS}{37}
\newcommand{\TTokCritiqueLowOSS}{945}
\newcommand{\TDebateLowOSS}{47}
\newcommand{\TTokDebateLowOSS}{1,040}
\newcommand{\TDirectLowGemma}{63}
\newcommand{\TTokDirectLowGemma}{520}
\newcommand{\TVoteLowGemma}{63}
\newcommand{\TTokVoteLowGemma}{1,490}
\newcommand{\TCritiqueLowGemma}{57}
\newcommand{\TTokCritiqueLowGemma}{884}
\newcommand{\TDebateLowGemma}{67}
\newcommand{\TTokDebateLowGemma}{1,378}
\newcommand{\ItemShareThirty}{47}
\newcommand{\ArmShareThirty}{3.2}
\newcommand{\ItemShareThirtyLo}{32}
\newcommand{\ItemShareThirtyHi}{59}
\newcommand{\ArmShareThirtyLo}{2.3}
\newcommand{\ArmShareThirtyHi}{9.2}
\newcommand{\NArmsThirty}{23}
\newcommand{\NItemsThirty}{30}
\newcommand{\ItemShareMain}{55}
\newcommand{\ArmShareMain}{0.5}
\newcommand{\ItemShareMainLo}{45}
\newcommand{\ItemShareMainHi}{63}
\newcommand{\ArmShareMainLo}{0.2}
\newcommand{\ArmShareMainHi}{2.4}
\newcommand{\NArmsMain}{8}
\newcommand{\NItemsMain}{83}

\newcommand{\vote}{\textsc{Vote}}
\newcommand{\critique}{\textsc{Critique}}
\newcommand{\debate}{\textsc{Debate}}
\newcommand{\direct}{\textsc{Direct}}

\title{Inference Can't Buy Knowledge: Agent Orchestration and Extended Reasoning\\Hit the Same Ceiling on Biology Research Tasks}
\author{Anonymous Authors}
\date{}

\begin{document}
\maketitle

\begin{abstract}
Agentic systems for biology usually report gains over a single model call while spending several calls.
We ask where such gains come from on objectively scored biology research tasks (LAB-Bench and LABBench2: protocol troubleshooting, sequence manipulation, database and literature questions).
At a matched budget of three calls, we compare independent sampling with majority vote, sequential self-critique, and a two-scientist debate with an adjudicator, across three open and API models on 150 pre-registered questions scored without an LLM judge, and we contrast these with simply letting frontier models reason longer.
Self-critique never beats majority voting (pooled $\DPooledCritiqueMinusVote$ points) and breaks at least as many answers as it repairs.
Debate helps only where the model cannot reason on its own: it adds \DDebateMinusVoteGemma{} points for Gemma with thinking disabled, but \DDebateMinusVoteOSS{} and \DDebateMinusVoteLite{} for models that reason, and no reliable gain for frontier models at high reasoning effort.
Raising reasoning effort rarely helps either: GPT-OSS generates \TokRatioOSS$\times$ more tokens at high effort for no gain, and across six model families accuracy varies little while spending varies by nearly two orders of magnitude.
The dominant factor is the question itself: across systems, question identity explains \ItemShareMain\% of the variance in correctness, while the choice of model, strategy and reasoning effort explains \ArmShareMain\%.
On these tasks, additional inference, however organised, does not substitute for missing knowledge.
\end{abstract}

\section{Introduction}
LLM agents now plan experiments, analyse data and generate hypotheses in chemistry and biology \citep{boiko2023coscientist,bran2023chemcrow,swanson2024virtuallab,huang2025biomni,gottweis2026coscientist,ghareeb2026robin}.
Their gains are typically reported against a single-call baseline, although the agent makes many calls.
Because repeated sampling with a vote already improves accuracy \citep{wang2022selfconsistency,li2024moreagents,brown2024monkeys}, such comparisons conflate \emph{what the interaction adds} with \emph{how much inference was spent} \citep{kapoor2024agentsmatters}.

We run a small, controlled test of this confound on biology research questions with deterministic answers.
Every orchestrated strategy gets exactly the budget of the voting baseline (three calls, same model, same prompts, same decoding).
We then ask a second question the agent literature rarely asks: does letting a strong model reason longer in a \emph{single} call reach the same point?

\paragraph{Findings.}
(1)~At matched budget, sequential critique never improves on majority voting, and multi-agent debate improves on it only for a model whose reasoning is switched off.
(2)~Critique never helps: it breaks at least as many correct answers as it repairs, and for models that reason it is most harmful on exact sequence manipulation.
(3)~Thinking longer is not a way out: higher reasoning effort helps only a model that otherwise does not reason at all, and then only up to the level its low-effort debate already reached; stacking orchestration on top adds nothing reliable.
(4)~Most strikingly, whether a system is right depends overwhelmingly on \emph{which question} it faces and barely on \emph{how} it spends inference: every system solves the same questions and fails the same questions (Fig.~\ref{fig:ceiling}).
The remaining errors look like missing knowledge, not missing deliberation.

\section{Setup}
\paragraph{Tasks.}
We use four LAB-Bench categories \citep{laurent2024labbench} (ProtocolQA, SeqQA, DbQA, LitQA2; multiple choice with shuffled options and no abstain option) and the deterministic subset of LABBench2 SeqQA2 \citep{labbench2} (open-ended sequence computations with unique answers, scored by the official validators, with input files injected into the prompt).
We exclude figure/table tasks (text-only models) and prompts over 6{,}000 characters.
With a fixed seed we draw 30 questions per category (150 total), plus 10 pilot questions that are excluded from all results.
No LLM judge is used anywhere.

\paragraph{Strategies (Fig.~\ref{fig:strategies} in App.~\ref{app:prompts}).}
\direct{} makes one call.
\vote{} makes three independent calls and takes the majority answer, breaking ties by self-reported confidence.
\critique{} solves, critiques the solution, then revises.
\debate{} has a molecular biologist and a skeptical reviewer solve independently, and a third call adjudicates between them.
\vote{}, \critique{} and \debate{} each use exactly three calls. Only visible outputs are passed between calls.

\paragraph{Models.}
The main study uses GPT-OSS-120B \citep{openai2025gptoss} (\texttt{reasoning\_effort=low}), Gemma~4~31B \citep{google2026gemma4} (thinking off, the only non-default setting the API accepts), and Gemini~3.5 Flash-Lite (thinking low), all through free API tiers.
Decoding uses temperature 0.6, fixed per-call seeds and a 2{,}048-token cap.
GPT-OSS and Gemma answer all 150 questions and Flash-Lite a fixed 75-question subset.
For the reasoning-effort analysis we additionally run GPT-6 Luna, GPT-6 Sol and GPT-5.6 Terra (all strategies, high effort; Luna also at low effort), GPT-OSS at high effort (all strategies, served by Ollama; a low-effort \direct{} run on the same host matches the main run's accuracy) and Flash-Lite (\direct{} at high thinking), all on the first 30 questions in the fixed run order (6 per category).

\paragraph{Pre-registration.}
Task IDs, prompts, decoding, hypotheses and the analysis script were committed before the main run.
The primary contrast is \debate{} $-$ \vote{}, tested with a paired bootstrap over questions (10{,}000 resamples) and an exact McNemar test.
Post-freeze deviations are logged in the repository. They are rate-limit workarounds, a second inference provider for 22 GPT-OSS questions (each question served entirely by one provider), and the exploratory reasoning-effort arm.

\section{Results}
\begin{table}[t]
\centering\small
\caption{Accuracy (\%) at a matched budget of three calls; \direct{} uses one call. $\Delta$ is \debate{} $-$ \vote{} in points with a paired bootstrap 95\% CI.}
\label{tab:main}
\begin{tabular}{lrrrrrc}
\toprule
Model & $n$ & \direct & \vote & \critique & \debate & $\Delta$ [95\% CI] \\
\midrule
GPT-OSS-120B & \NOSS & \AccDirectOSS & \AccVoteOSS & \AccCritiqueOSS & \AccDebateOSS & \DDebateMinusVoteOSS{} [\DDebateMinusVoteOSSLo, \DDebateMinusVoteOSSHi] \\
Gemma 4 31B & \NGemma & \AccDirectGemma & \AccVoteGemma & \AccCritiqueGemma & \AccDebateGemma & \DDebateMinusVoteGemma{} [\DDebateMinusVoteGemmaLo, \DDebateMinusVoteGemmaHi] \\
Flash-Lite & \NLite & \AccDirectLite & \AccVoteLite & \AccCritiqueLite & \AccDebateLite & \DDebateMinusVoteLite{} [\DDebateMinusVoteLiteLo, \DDebateMinusVoteLiteHi] \\
\bottomrule
\end{tabular}
\end{table}

\paragraph{Orchestration beats voting only when the model cannot think.}
\debate{} improves on \vote{} for Gemma, whose thinking is disabled ($\DDebateMinusVoteGemma$ points, 95\% CI [$\DDebateMinusVoteGemmaLo$, $\DDebateMinusVoteGemmaHi$], McNemar $p=\PDebateMinusVoteGemma$), and the gain survives excluding every question with a truncated call.
For the two models that reason, the interval includes zero ($p=\PDebateMinusVoteOSS$ and $\PDebateMinusVoteLite$; for GPT-OSS the estimate is unchanged at $\DDebateMinusVoteOSSNoGroq$ when the \the\numexpr150-\NOSSNoGroq\relax{} questions served by the second provider are excluded).
Pooled over all \NPooled{} model--question pairs, \debate{} $-$ \vote{} is $\DPooledDebateMinusVote$ [$\DPooledDebateMinusVoteLo$, $\DPooledDebateMinusVoteHi$] points and \critique{} $-$ \vote{} is $\DPooledCritiqueMinusVote$ [$\DPooledCritiqueMinusVoteLo$, $\DPooledCritiqueMinusVoteHi$].
\critique{} is never better than \vote{} ($\DCritiqueMinusVoteOSS$, $\DCritiqueMinusVoteGemma$, $\DCritiqueMinusVoteLite$ points).
Voting itself barely improves on a single call ($\DVoteMinusDirectOSS$, $\DVoteMinusDirectGemma$, $\DVoteMinusDirectLite$ points), even though the three samples disagree on \DisFracOSS\%, \DisFracGemma\% and \DisFracLite\% of questions.

\paragraph{Critique breaks at least as often as it repairs.}
Relative to the vote, \critique{} rescues \RescueCritiqueOSS\%, \RescueCritiqueGemma\% and \RescueCritiqueLite\% of wrong answers but corrupts \CorruptCritiqueOSS\%, \CorruptCritiqueGemma\% and \CorruptCritiqueLite\% of correct ones.
For the two models that reason, the largest losses are on exact sequence manipulation (SeqQA: \CritMinusVoteSeqQAOSS{} and \CritMinusVoteSeqQALite{} points vs.\ \vote{} for GPT-OSS and Flash-Lite; \CritMinusVoteSeqQAGemma{} for Gemma), where a fluent but wrong critique talks the model out of a correct computation.

\begin{figure}[t]
\centering
\figifexists[width=\linewidth]{figures/fig_think_vs_talk.pdf}
\caption{Accuracy vs.\ generated tokens per question on the same 30 questions. Circles are low reasoning and stars are high reasoning; colour marks the strategy. Across nearly two orders of magnitude in generated tokens, all ways of spending inference land on a narrow per-model plateau.}
\label{fig:tvt}
\end{figure}

\paragraph{Thinking longer mostly does not help either.}
On the same 30 questions (Fig.~\ref{fig:tvt}), raising GPT-OSS to high reasoning effort multiplies generated tokens by \TokRatioOSS$\times$ (\TTokDirectLowOSS{} to \TTokDirectHighOSS{} per question) while accuracy stays flat (\TDirectLowOSS\% to \TDirectHighOSS\%). For Flash-Lite the result is the same (\TDirectLowLite\% to \TDirectHighLite\%, \TokRatioLite$\times$ tokens).
The one model that gains is GPT-6 Luna, which at low effort uses no reasoning tokens (\TDirectLowLuna\% to \TDirectHighLuna\%). Its high-effort single call lands exactly where its low-effort \debate{} already was (\TDebateLowLuna\%), consistent with debate acting as a substitute for reasoning.
Orchestration on top of high effort adds nothing reliable. GPT-6 Sol scores \TDirectHighSol\% with all four strategies, and GPT-5.6 Terra ranges from \TCritiqueHighTerra\% (\critique) to \TDebateHighTerra\% (\debate) with no paired difference significant.
Within each model, accuracy moves by a few questions while generated tokens vary by up to two orders of magnitude.

\begin{figure}[t]
\centering
\figifexists[width=\linewidth]{figures/fig_ceiling_main.pdf}
\caption{Correctness of every model $\times$ strategy system (rows) on every question answered by all of them (columns, sorted by the fraction of systems correct). The staircase shows that systems solve and fail the \emph{same} questions and differ mainly in a shared boundary. Question identity explains \ItemShareMain\% [\ItemShareMainLo, \ItemShareMainHi] of the variance in correctness; the system explains \ArmShareMain\% [\ArmShareMainLo, \ArmShareMainHi].}
\label{fig:ceiling}
\end{figure}

\paragraph{The question, not the method, determines the outcome.}
A two-way decomposition of the 0/1 correctness matrix (systems $\times$ questions) attributes \ItemShareMain\% of the variance to the question and \ArmShareMain\% to the system for the main models (Fig.~\ref{fig:ceiling}).
Across all \NArmsThirty{} systems on the shared 30 questions, including frontier models at high reasoning, the split is \ItemShareThirty\% vs.\ \ArmShareThirty\%.
Of the \NUnsolved{} questions that no main system solves, \NUnsolvedRecall{} require literature or database recall (LitQA2, DbQA) and \NUnsolvedSeqQAII{} are open-ended sequence computations (SeqQA2), failures that point to missing retrieval and code execution rather than missing deliberation.

\section{Discussion}
On objectively scored biology research tasks, the additional calls in agentic pipelines buy little beyond what independent sampling or a longer single call buys.
Debate is the exception, and only when the model is not allowed to reason; once reasoning is on, its gain disappears, and critique can hurt when a confident critic overrides a correct computation.
Longer reasoning is not a reliable substitute either: it helped only a model that otherwise did not reason.
The practical implication for agent builders is to measure against a compute-matched vote \citep{kapoor2024agentsmatters} and to invest the budget in \emph{knowledge access} (retrieval, databases, code execution) rather than in more elaborate conversations between copies of the same model.
Our results agree with compute-matched analyses of debate on general reasoning benchmarks \citep{wang2024rethinking,smit2024mad,choi2025debatevsvote,wu2025candebate} and with evidence that intrinsic self-correction is unreliable \citep{huang2023cannotcorrect}. We extend them to biology, to sequential critique, and to the trade-off between orchestration and reasoning effort.

\paragraph{Limitations.}
Models had no tools or retrieval, so DbQA and LitQA2 measure recall; orchestration with tools may behave differently.
The budget is fixed at three calls.
Frontier models were run through the Codex CLI, which adds its own system prompt and does not expose temperature or seeds; they cover 30 questions, so their comparisons are exploratory.
Gemma could only run with thinking off.
One prompt set per strategy was used, and 150 questions give limited power for effects under ${\sim}10$ points.

\bibliographystyle{plainnat}
\bibliography{refs}

\appendix
\section{Prompts}\label{app:prompts}
All prompts are single user messages. \texttt{\{QUESTION\}} is the question (with protocol, options or injected files) and \texttt{\{ANSWER INSTRUCTION\}} is the format line for multiple choice or the official SeqQA2 suffix.
\refstepcounter{figure}\label{fig:strategies}
\paragraph{Direct / Independent$\times$3 / Critique call 1 (solve)}
\begin{footnotesize}\begin{verbatim}
Solve the following biology research question.

{QUESTION}

{ANSWER INSTRUCTION}

Respond in exactly this format and keep it short:
<answer>YOUR ANSWER</answer>
CONFIDENCE: <probability between 0 and 1 that your answer is correct>
JUSTIFICATION: <at most two sentences>
\end{verbatim}\end{footnotesize}

\paragraph{Critique call 2 (critique)}
\begin{footnotesize}\begin{verbatim}
Below are a biology research question and a proposed solution. Act as a critical reviewer: identify any
scientific, computational or reasoning error in the solution. If you think the answer is wrong, say which
answer you believe is correct and why. If you find no error, say so.

{QUESTION}

Proposed solution:
{CALL 1 OUTPUT}

Respond with at most four sentences, starting with "CRITIQUE:".
\end{verbatim}\end{footnotesize}

\paragraph{Critique call 3 (revise)}
\begin{footnotesize}\begin{verbatim}
Below are a biology research question, a first attempt at solving it, and a reviewer's critique of that
attempt. Either may contain mistakes. Considering both, produce the final answer.

{QUESTION}

First attempt:
{CALL 1 OUTPUT}

Reviewer's critique:
{CALL 2 OUTPUT}

{ANSWER INSTRUCTION}

Respond in exactly this format and keep it short:
<answer>YOUR ANSWER</answer>
CONFIDENCE: <probability between 0 and 1 that your answer is correct>
JUSTIFICATION: <at most two sentences>
\end{verbatim}\end{footnotesize}

\paragraph{Debate call 1 (researcher)}
\begin{footnotesize}\begin{verbatim}
You are a molecular biology researcher. Solve the following biology research question independently.

{QUESTION}

{ANSWER INSTRUCTION}

Respond in exactly this format and keep it short:
<answer>YOUR ANSWER</answer>
CONFIDENCE: <probability between 0 and 1 that your answer is correct>
JUSTIFICATION: <at most two sentences>
\end{verbatim}\end{footnotesize}

\paragraph{Debate call 2 (skeptic)}
\begin{footnotesize}\begin{verbatim}
You are a skeptical scientific reviewer. Solve the following biology research question independently, looking
especially for subtle experimental, biological or computational errors and traps.

{QUESTION}

{ANSWER INSTRUCTION}

Respond in exactly this format and keep it short:
<answer>YOUR ANSWER</answer>
CONFIDENCE: <probability between 0 and 1 that your answer is correct>
JUSTIFICATION: <at most two sentences>
\end{verbatim}\end{footnotesize}

\paragraph{Debate call 3 (adjudicator)}
\begin{footnotesize}\begin{verbatim}
Two scientists independently analysed the biology research question below. Adjudicate between their analyses
(either or both may be wrong) and submit the final answer.

{QUESTION}

Analysis 1 (molecular biology researcher):
{CALL 1 OUTPUT}

Analysis 2 (skeptical scientific reviewer):
{CALL 2 OUTPUT}

{ANSWER INSTRUCTION}

Respond in exactly this format and keep it short:
<answer>YOUR ANSWER</answer>
CONFIDENCE: <probability between 0 and 1 that your answer is correct>
JUSTIFICATION: <at most two sentences>
\end{verbatim}\end{footnotesize}


\section{Additional results}
\begin{table}[h]
\centering
\small
\begin{tabular}{lccccccccccccccc}
\toprule
Category & \multicolumn{5}{c}{Gemma} & \multicolumn{5}{c}{Lite} & \multicolumn{5}{c}{OSS} \\
 & D & I & C & M & Disagree & D & I & C & M & Disagree & D & I & C & M & Disagree \\
\midrule
DbQA & 56.7 & 53.3 & 40.0 & 56.7 & 33.3 & 40.0 & 40.0 & 40.0 & 33.3 & 26.7 & 39.1 & 43.5 & 47.8 & 34.8 & 34.8 \\
LitQA2 & 33.3 & 30.0 & 33.3 & 36.7 & 46.7 & 60.0 & 53.3 & 53.3 & 46.7 & 26.7 & 21.7 & 26.1 & 43.5 & 39.1 & 47.8 \\
ProtocolQA & 56.7 & 56.7 & 60.0 & 70.0 & 20.0 & 66.7 & 80.0 & 86.7 & 73.3 & 26.7 & 62.5 & 50.0 & 45.8 & 50.0 & 33.3 \\
SeqQA & 56.7 & 56.7 & 66.7 & 73.3 & 46.7 & 73.3 & 73.3 & 53.3 & 66.7 & 13.3 & 62.5 & 66.7 & 41.7 & 62.5 & 33.3 \\
SeqQA2 & 40.0 & 46.7 & 36.7 & 50.0 & 63.3 & 66.7 & 66.7 & 46.7 & 66.7 & 40.0 & 45.8 & 50.0 & 41.7 & 58.3 & 41.7 \\
\bottomrule
\end{tabular}
\caption{Per-category accuracy (D=Direct, I=Independent×3, C=Critique, M=Multi). Numbers are percentages.}
\label{tab:by-category}
\end{table}

\end{document}
